\Needspace{6\baselineskip}
\begin{origexamplebox}{例题 9.11【貌似简谐振动】}
半径为 $R$ 的圆柱形区域内有垂直纸面向里的匀强磁场
$B=B_0\sin\omega t$。在距离圆心 $h$ 的一条弦上有光滑直轨，带电质点 $m,+q$ 被约束在轨道中，初始位于弦中点且静止。

(1) 求粒子的运动规律和任意时刻的位置；

(2) 求保证粒子不离开磁场区域的参数条件。

% Pure source-question geometry: no derived forces, currents or solution curves.
%
\begin{center}\begin{tikzpicture}[scale=0.8,>=Latex]
  \draw[Ink,thick] (0,0) circle (2.5);
  \draw[Ink,very thick] ({-sqrt(2.5*2.5-1.25*1.25)},-1.25)--({sqrt(2.5*2.5-1.25*1.25)},-1.25);
  \fill[DeepBlue] (0,-1.25) circle (3pt) node[below] {$m,+q$};
  \draw[<->] (0,0)--node[right] {$h$}(0,-1.25);
  \draw[axis] (-2.2,-1.75)--(2.4,-1.75) node[right] {$x$};
  \node[TextGray] at (-1.4,1.0) {$\times$}; \node[TextGray] at (0,1.0) {$\times$}; \node[TextGray] at (1.4,1.0) {$\times$};
  \node at (0,1.65) {$B_0\sin\omega t$};
  \end{tikzpicture}\end{center}
\figcap{正电荷位于弦中点，圆形区域内磁场按给定规律变化。}\label{fig:em24-question-11}
\end{origexamplebox}

\begin{xesolution}
取轨道方向向右为 $x$ 轴，原点在弦中点。对于圆柱内的匀强时变磁场，感生电场可写为
\begin{equation}
  \vect E=-\frac12\dot B_z\,\uvec z\times\vect r.
\end{equation}
磁场指向纸内，在圆心下方距离 $h$ 的弦上，电场沿轨道的分量与 $x$ 无关：
\begin{equation}
  E_x=\frac h2\dot B=\frac{hB_0\omega}{2}\cos\omega t.
\end{equation}
磁力与轨道方向垂直，由轨道的支持力平衡，因而只有电场力改变 $x$ 向运动：
\begin{equation}
  m\ddot x=\frac{qhB_0\omega}{2}\cos\omega t.
\end{equation}
由 $x(0)=0,\dot x(0)=0$ 得
\begin{equation}
  \boxed{\dot x(t)=\frac{qhB_0}{2m}\sin\omega t},
  \qquad
  \boxed{x(t)=\frac{qhB_0}{2m\omega}(1-\cos\omega t)}.
\end{equation}
这是绕平衡位置 $x_c=qhB_0/(2m\omega)$ 的简谐形运动，但驱动来自时变磁场，并非恢复力。弦的半长为
$s=\sqrt{R^2-h^2}$，而 $|x|$ 的最大值为 $|q|hB_0/(m\omega)$，故条件为
\begin{equation}
  \boxed{\frac{|q|hB_0}{m\omega}\le\sqrt{R^2-h^2}}.
\end{equation}

图~\ref{ch24:fig:review-chordparticlediagram-1}中，轨道在圆心下方h处，端点由同一个圆与直线的交点确定；粒子位移不能越过弦端。
\par\noindent\begin{minipage}{\linewidth}
\chordParticleDiagram
\figcap{弦上粒子的约束与磁场边界}\label{ch24:fig:review-chordparticlediagram-1}
\end{minipage}\par\medskip
\end{xesolution}

\Needspace{6\baselineskip}
\begin{origexamplebox}{例题 9.12【感生电动势】}
半径为 $R$ 的圆柱区域内有垂直纸面向里的匀强磁场
$B=B_0+kt$。距圆心 $h$ 处有一条直导线 $AD$，与圆的两个交点为 $A,C$，$D$ 位于 $C$ 右侧。记 $CD=\ell$，求沿 $A\to C$、$C\to D$ 两段的感生电动势，并讨论 $D$ 趋于无穷远的特例。

% Pure source-question geometry: no derived forces, currents or solution curves.
%
\begin{center}\begin{tikzpicture}[scale=0.82,>=Latex]
  \draw[Ink,thick] (0,0) circle (2.5);
  \draw[Ink,very thick] (-3.2,-1.2)--(4.3,-1.2);
  \coordinate (A) at ({-sqrt(2.5*2.5-1.2*1.2)},-1.2); \coordinate (C) at ({sqrt(2.5*2.5-1.2*1.2)},-1.2); \coordinate (D) at (4.1,-1.2);
  \fill[DeepBlue] (A) circle (2.5pt) node[below] {$A$}; \fill[DeepBlue] (C) circle (2.5pt) node[below] {$C$}; \fill[Ink] (D) circle (2.5pt) node[below] {$D$};
  \draw[<->] ({sqrt(2.5*2.5-1.2*1.2)},-2)--(4.1,-2) node[midway,below] {$\ell$};
  \draw[<->] (0,0)--node[right] {$h$}(0,-1.2);
  \node[TextGray] at (-1.1,.2) {$\dot B=k$};
\end{tikzpicture}\end{center}
\figcap{弦、观察端点和外延直线；感生电场分布待求。}\label{fig:em24-question-12}
\end{origexamplebox}

\begin{xesolution}
先取 $0<h<R$，令 $s=\sqrt{R^2-h^2}$，以垂足为坐标原点，导线向右为积分正向；则 $x_A=-s$、$x_C=s$、$x_D=s+\ell$。磁场向里增大时，感生电场从读者视角看为逆时针，因而在圆心下方的导线上其分量向右。

圆内电场 $E_\varphi=kr/2$。在水平线 $y=-h$ 上，其 $x$ 分量为常量 $kh/2$，故
\begin{equation}
  \boxed{\mathcal E_{AC}=\int_{-s}^{s}\frac{kh}{2}\dd x
  =kh\sqrt{R^2-h^2}}.
\end{equation}
圆外电场 $E_\varphi=kR^2/(2r)$，而水平分量为
\begin{equation}
  E_x=\frac{kR^2h}{2(x^2+h^2)}.
\end{equation}
因此对于有限端点 $D$，
\begin{align}
  \mathcal E_{CD}
  &=\frac{kR^2h}{2}\int_s^{s+\ell}\frac{\dd x}{x^2+h^2}\\
  &=\boxed{\frac{kR^2}{2}\left[\arctan\frac{s+\ell}{h}-\arctan\frac{s}{h}\right]}.
\end{align}
只有再令 $\ell\to\infty$，才得到
\[
 \boxed{\mathcal E_{C\infty}=\frac{kR^2}{2}\left(\frac\pi2-\arctan\frac{s}{h}\right)
 =\frac{kR^2}{2}\arcsin\frac hR}.
\]
当 $\ell=0$ 时第二段电动势为零；当 $h=0$ 时两段上的电场沿线分量都为零，可由几何直接判断。
若改变积分方向，电动势整体变号。
后一段虽然完全位于无磁场区域，感生电场却并不为零；这正是涡旋电场与“局部磁力”图像的区别。

图~\ref{ch24:fig:review-chordemfdiagram-1}中，A、C为圆与导线的交点，D为有限观察端点；圆外仍需积分涡旋电场。
\par\noindent\begin{minipage}{\linewidth}
\chordEmfDiagram
\figcap{跨越磁场边界的直导线}\label{ch24:fig:review-chordemfdiagram-1}
\end{minipage}\par\medskip
\end{xesolution}

\Needspace{6\baselineskip}
\begin{origexamplebox}{例题 9.13【这是测的啥】}
半径为 $R$ 的圆柱区域内有垂直纸面向里的匀强磁场 $B=B_0+kt$。边缘处放置一根不均匀的圆形电阻线，两端点为 $A,C$。
劣弧 $AC$ 的圆心角为 $2\pi/3$、电阻为 $r_0/4$；优弧 $AC$ 的圆心角为 $4\pi/3$、电阻为 $r_0/3$。电压表 $V_1,V_2$ 分别紧贴劣弧和优弧连接 $A,C$；
$V_3$ 的两根导线从 $A,C$ 沿径向通到圆心。求三表示数并解释它们测量的物理量。
\par\noindent\begin{minipage}{\linewidth}
\circularVoltmetersDiagram
\figcap{圆形电阻线与三种测量路径}\label{ch24:fig:review-circularvoltmetersdiagram-1}
\end{minipage}\par\medskip
\end{origexamplebox}

\begin{xesolution}
总感应电动势的大小为
\begin{equation}
  \mathcal E_0=k\pi R^2.
\end{equation}
电阻线中的电流为
\begin{equation}
  I=\frac{\mathcal E_0}{r_0/4+r_0/3}
  =\frac{12\mathcal E_0}{7r_0}.
\end{equation}
$V_1,V_2$ 的导线各自紧贴被测电阻弧，电压表与电阻弧围成的面积可趋于零，所以它们读出的是相应电阻上的 Ohm 压降：
\begin{equation}
  \boxed{V_1=I\frac{r_0}{4}=\frac37\mathcal E_0},
  \qquad
  \boxed{V_2=I\frac{r_0}{3}=\frac47\mathcal E_0}.
\end{equation}
$V_3$ 的两根引线都沿径向，引线上 $\vect E_{\rm ind}\cdot\dd\vect\ell=0$。劣弧和两根径向线围成的扇形面积为 $\pi R^2/3$，其感应电动势为 $\mathcal E_0/3$。因此
\begin{equation}
  \boxed{V_3=\left|I\frac{r_0}{4}-\frac{\mathcal E_0}{3}\right|
  =\frac{2}{21}\mathcal E_0}.
\end{equation}
$V_3$ 测得的是由该引线路径定义的端点读数，它不等于任意一段电阻压降。时变磁场中 $\oint\vect E\cdot\dd\vect\ell\ne0$，“两点间电压”会依赖电压表引线的路径。

图~\ref{ch24:fig:review-circularvoltmetersdiagram-1}中，V1、V2的引线分别紧贴两段弧，V3沿两条半径；不同闭合路径对应不同包围面积。

\end{xesolution}

\Needspace{6\baselineskip}
\begin{origexamplebox}{例题 9.14【这也可以测】}
边长为 $a$ 的正方形导线框放在垂直纸面向里的匀强磁场 $B=B_0+kt$ 中。左、右两边电阻分别为 $R_1,R_2$，上、下两边电阻忽略。
电压表 $V_1,V_2,V_3,V_4$ 分别紧贴上、右、下、左四边连接。求各表示数并解释其物理意义。
\par\noindent\begin{minipage}{\linewidth}
\squareVoltmetersDiagram
\figcap{四只电压表的真实连接点}\label{ch24:fig:review-squarevoltmetersdiagram-1}
\end{minipage}\par\medskip
\end{origexamplebox}

\begin{xesolution}
线框的感应电动势大小为 $\mathcal E_0=ka^2$，电流大小为
\begin{equation}
  I=\frac{ka^2}{R_1+R_2}.
\end{equation}
紧贴某一边的电压表与该边围成的面积趋于零，因此其读数就是该边的 Ohm 压降。上下两边无电阻，故
\begin{equation}
  \boxed{V_1=0},
  \qquad
  \boxed{V_2=\frac{R_2}{R_1+R_2}ka^2},
\end{equation}
\begin{equation}
  \boxed{V_3=0},
  \qquad
  \boxed{V_4=\frac{R_1}{R_1+R_2}ka^2}.
\end{equation}
虽然 $V_1$ 和 $V_3$ 都连在左右两侧导体之间，但它们的引线路径不同。在涡旋电场中，不能脱离测量路径把“两点电压”当成唯一的标量。

图~\ref{ch24:fig:review-squarevoltmetersdiagram-1}中，所有引线都接到顶点，表与被测边的间距仅为辨认而放大，计算取其所围附加面积趋零。

\end{xesolution}

\Needspace{6\baselineskip}
\begin{origexamplebox}{例题 9.15【奇妙的田字格】}
一个外框边长为 $2a$ 的“田”字格电阻网络放在垂直纸面向里的匀强磁场 $B=B_0+kt$ 中。每根导线的单位长度电阻为 $r/a$，因而每段长为 $a$ 的小边电阻为 $r$。求各支路电流。

% Pure source-question geometry: no derived forces, currents or solution curves.
%
\begin{center}\begin{tikzpicture}[scale=.85,>=Latex]
 \draw[Ink,thick](-2,-2)grid[step=2](2,2);
 \node at(-1,1){$\otimes$};\node at(1,1){$\otimes$};\node at(-1,-1){$\otimes$};\node at(1,-1){$\otimes$};
 \draw[<->](-2,-2.5)--(2,-2.5)node[midway,below]{$2a$};\node[right]at(2.6,.2){$\dot B=k$};
\end{tikzpicture}\end{center}
\figcap{四个方格中的垂直磁场按相同变化率改变。}\label{fig:em24-question-15}
\end{origexamplebox}

\begin{xesolution}
四个边长为 $a$ 的小方格所围面积相同，每个网孔的感应电动势都是
\begin{equation}
  \mathcal E_c=ka^2,
\end{equation}
且方向都为逆时针。由四重旋转对称性，四个网孔电流必然相等，记为 $i$。每一条内部十字支路都是两个相邻网孔的公用边，两个等大网孔电流在其上方向相反，故内部支路的实际电流为零。

对任一小网孔写 Kirchhoff 方程，只有位于外边界的两段电阻产生压降：
\begin{equation}
  2ri=ka^2.
\end{equation}
因此
\begin{equation}
  \boxed{i=\frac{ka^2}{2r}}.
\end{equation}
外周八段小边均通有大小 $ka^2/(2r)$ 的逆时针电流，中央十字的四段电流均为零。

图~\ref{ch24:fig:review-gridcurrentdiagram-1}中，向纸内磁场增强时，外周逆时针电流产生向纸外的反抗磁场，内部支路电流为零。
\par\noindent\begin{minipage}{\linewidth}
\gridCurrentDiagram
\figcap{田字格的外周环流}\label{ch24:fig:review-gridcurrentdiagram-1}
\end{minipage}\par\medskip
\end{xesolution}

\Needspace{6\baselineskip}
\begin{origexamplebox}{例题 9.16【万能的 Kirchhoff】}
证明：任意复杂的、仅含电阻和理想电源的线性网络放入匀强时变磁场 $B=B_0+kt$ 中，只要所有支路参数和各独立回路所围面积均已知，就能求出所有支路电流；并证明此时仍满足电源叠加定理。
\end{origexamplebox}

\begin{xesolution}
为网络选取 $n$ 个独立网孔电流，组成列向量 $\vect i$。记 $\vect R$ 为由支路电阻组成的网孔电阻矩阵，$\vect e$ 为理想电源在各网孔中的代数和，
$\vect S=(S_1,\ldots,S_n)^{\mathsf T}$ 为有符号面积列向量。Faraday--Kirchhoff 方程可写成
\begin{equation}
  \boxed{\vect R\vect i=\vect e-k\vect S}.
\end{equation}
对于非平面网络，“网孔”应改用一般的独立回路，$S_j$ 为其有向面积在磁场方向的投影。
设支路--回路矩阵为 $\vect A$，支路电阻矩阵为 $\vect D=\operatorname{diag}(R_e)$，
则 $\vect R=\vect A^{\mathsf T}\vect D\vect A$，且
$\vect i^{\mathsf T}\vect R\vect i=\sum_eR_eI_e^2\ge0$。
这只能无条件保证半正定。若每个非零独立环流都经过正电阻，才有严格正定及可逆性。所以在此条件下
\begin{equation}
  \boxed{\vect i=\vect R^{-1}\vect e-k\vect R^{-1}\vect S}.
\end{equation}
再用网孔--支路关联矩阵即可恢复所有支路电流。上式的第一项是普通电源单独作用的响应，第二项是时变磁场等效电源单独作用的响应，两者线性相加，故叠加定理仍然成立。
若存在纯理想电源或零电阻回路，必须先检查约束：回路电动势不相容时无有限稳态解，
相容时也可能不能唯一确定电流。例如两只等电压理想电源并联，其分担电流不由理想模型唯一决定。
因此题设的“任意”须理解为非退化、约束相容的线性网络；退化情况不能直接写 $\vect R^{-1}$。
\end{xesolution}

\Needspace{6\baselineskip}
\begin{origexamplebox}{例题 9.17【自感线圈】}
一个理想密绕长螺线管的长度为 $l$、半径为 $R$、单位长度匝数为 $n$。

(1) 求其自感；

(2) 若把它看成左右两部分，整体自感是否等于两部分自感之和？

% Pure source-question geometry: no derived forces, currents or solution curves.
%
\begin{center}\begin{tikzpicture}[scale=.75,>=Latex]
 \foreach \x in {-3.6,-3.2,...,3.6}{\draw[Ink](\x,0)ellipse(.18 and 1.0);}
 \draw[guide,dashed](0,-1.4)--(0,1.5);
 \draw[<->](-3.8,-1.5)--(0,-1.5)node[midway,below]{$l_1$};
 \draw[<->](0,-1.5)--(3.8,-1.5)node[midway,below]{$l_2$};
 \end{tikzpicture}\end{center}
\figcap{沿同一长螺线管把绕组分为两段，位置保持不变。}\label{fig:em24-question-17}
\end{origexamplebox}

\begin{xesolution}
螺线管内部磁场为 $B=\mu_0nI$，截面积 $A=\pi R^2$，总匝数 $N=nl$。磁链为
\begin{equation}
  \Lambda=N\Phi=(nl)(\mu_0nI)(\pi R^2).
\end{equation}
所以
\begin{equation}
  \boxed{L=\frac\Lambda I=\mu_0n^2\pi R^2l}.
\end{equation}
若左右两段长度为 $l_1,l_2$，且 $l_1+l_2=l$，先保持两段原有位置和绕向。
同一电流 $I$ 依次通过两段，磁能为
\[
 W=\frac12L_1I^2+MI^2+\frac12L_2I^2
   =\frac12(L_1+L_2+2M)I^2,
\]
所以一般关系是 $L=L_1+L_2+2M$，而不是无条件简单相加。
若两个子段都满足 $l_1,l_2\gg R$，并在长管主阶近似中忽略分界附近的交链端场，则
\begin{equation}
  L_1+L_2\simeq\mu_0n^2\pi R^2(l_1+l_2)\simeq L.
\end{equation}
因此“等于”只指这一近似的主阶结果。整体长不代表任意切出的短段也长；如果一段长度与半径相当，
不能对它继续使用长管公式。单纯把一根线圈在账面上分成两段不会改变其几何或磁场，互感项仍需计入；
实际切开并移远则改变了互感，是另一个装置。

图~\ref{ch24:fig:review-solenoidsplitdiagram-1}中，分界只是对同一线圈的划分；保持位置时总自感还包括两段互感的两倍。
\par\noindent\begin{minipage}{\linewidth}
\solenoidSplitDiagram
\figcap{同向串联的两段线圈}\label{ch24:fig:review-solenoidsplitdiagram-1}
\end{minipage}\par\medskip
\end{xesolution}

\Needspace{6\baselineskip}
\begin{origexamplebox}{例题 9.18【互感线圈】}
两个单匝线圈放置在任意位置，分别通入电流 $I_1,I_2$。证明：1 号线圈产生的磁场穿过 2 号线圈的磁通量与 $I_1$ 之比，等于 2 号线圈产生的磁场穿过 1 号线圈的磁通量与 $I_2$ 之比，即 $M_{12}=M_{21}$。
% Question contract: only the two loops and prescribed current reference directions.
\begin{center}\begin{tikzpicture}[scale=.82,>=Latex]
 \draw[Ink,thick](-2.5,0)circle(1.2);\draw[Ink,thick](2.5,.4)circle(1);
 \draw[vector](-1.3,0)arc[start angle=0,end angle=65,radius=1.2];
 \draw[vector](3.5,.4)arc[start angle=0,end angle=65,radius=1];
 \node at(-2.5,0){线圈1};\node at(2.5,.4){线圈2};
 \node[right]at(-1.35,.95){$I_1$};\node[right]at(3.45,1.2){$I_2$};
\end{tikzpicture}\end{center}
\figcap{两个单匝线圈及各自的电流参考方向。}\label{fig:em24-question-18}

\end{origexamplebox}

\begin{xesolution}
\par\noindent\begin{minipage}{\linewidth}
\mutualLoopsProofDiagram
\figcap{互感双重积分的源点和场点}\label{ch24:fig:review-mutualloopsproofdiagram-1}
\end{minipage}\par\medskip

使用稳恒线电流的矢势
\begin{equation}
  \vect A_1(\vect r_2)=\frac{\mu_0I_1}{4\pi}
  \oint_{C_1}\frac{\dd\vect\ell_1}{r_{12}}.
\end{equation}
由 Stokes 定理，1 号线圈在 2 号线圈中建立的磁通为
\begin{align}
  \Phi_{2\leftarrow1}
  &=\oint_{C_2}\vect A_1\cdot\dd\vect\ell_2\\
  &=\frac{\mu_0I_1}{4\pi}\oint_{C_2}\!\oint_{C_1}
  \frac{\dd\vect\ell_1\cdot\dd\vect\ell_2}{r_{12}}.
\end{align}
因此
\begin{equation}
  M_{12}=\frac{\Phi_{2\leftarrow1}}{I_1}
  =\frac{\mu_0}{4\pi}\oint_{C_1}\!\oint_{C_2}
  \frac{\dd\vect\ell_1\cdot\dd\vect\ell_2}{r_{12}}.
\end{equation}
交换积分次序及下标 $1,2$，点积和距离 $r_{12}$ 均不变，所以
\begin{equation}
  \boxed{M_{12}=M_{21}}.
\end{equation}
两回路的形状、大小和相对位置可以任意；以上证明采用真空中的稳恒线电流近似。
普通线性互易介质也有互易性，但不能把真空积分核无条件照搬到任意介质。

图~\ref{ch24:fig:review-mutualloopsproofdiagram-1}中，两个线元落在各自回路上并取切向，交换下标不改变点积及两点距离。

\end{xesolution}

\Needspace{6\baselineskip}
\begin{origexamplebox}{例题 9.19【RLC 回路】}
电阻 $R$、电感 $L$、电容 $C$ 串联成回路。初始电流为零，电容初电压为 $U$。$t=0$ 闭合开关，求任意时刻的放电电流。
\par\noindent\begin{minipage}{\linewidth}
\rlcDiagram{0}
\figcap{RLC回路及初始电压参考极性}\label{ch24:fig:review-rlcdiagram-1}
\end{minipage}\par\medskip
\end{origexamplebox}

\begin{xesolution}
记电容器带电量为 $q$，并把放电方向的电流取为 $I=-\dot q$。则
\begin{equation}
  L\ddot q+R\dot q+\frac qC=0,
  \qquad q(0)=CU,
  \qquad \dot q(0)=0.
\end{equation}
特征根为
\begin{equation}
  s_\pm=-\frac{R}{2L}\pm
  \sqrt{\left(\frac{R}{2L}\right)^2-\frac1{LC}}.
\end{equation}
对 $s_+\ne s_-$，电流的统一表达式为
\begin{equation}
  \boxed{I(t)=\frac UL\,
  \frac{\e^{s_+t}-\e^{s_-t}}{s_+-s_-}}.
\end{equation}
当 $R<2\sqrt{L/C}$ 时，记
\begin{equation}
  \alpha=\frac{R}{2L},
  \qquad
  \omega_d=\sqrt{\frac1{LC}-\alpha^2},
\end{equation}
则欠阻尼电流为
\begin{equation}
  \boxed{I(t)=\frac{U}{L\omega_d}\e^{-\alpha t}\sin\omega_dt}.
\end{equation}
临界阻尼 $R=2\sqrt{L/C}$ 时，
\begin{equation}
  \boxed{I(t)=\frac{U}{L}\,t\e^{-Rt/(2L)}}.
\end{equation}
过阻尼情形直接使用实根形式。电流的初值为零是电感电流不能突变的结果，而非“电容还没开始放电”。

图~\ref{ch24:fig:review-rlcdiagram-1}中，图示开关已闭合，C为等长双极板，R、L使用独立元件符号；电流正向与初始电压极性共同决定方程。

\end{xesolution}

\Needspace{6\baselineskip}
\begin{origexamplebox}{例题 9.20【电源凑热闹】}
电阻 $R$、电感 $L$、电容 $C$ 与恒定电源 $\mathcal E$ 串联成回路。初始电流为零，电容初电压为 $U$。$t=0$ 闭合开关，求任意时刻的电流。
\par\noindent\begin{minipage}{\linewidth}
\rlcDiagram{1}
\figcap{RLC回路及初始电压参考极性}\label{ch24:fig:review-rlcdiagram-2}
\end{minipage}\par\medskip
\end{origexamplebox}

\begin{xesolution}
取电源充电方向为电流正方向，令 $q$ 为电容正极板电量，$I=\dot q$。回路方程为
\begin{equation}
  L\ddot q+R\dot q+\frac qC=\mathcal E,
  \qquad q(0)=CU,
  \qquad \dot q(0)=0.
\end{equation}
把 $q=C\mathcal E+y$ 代入，$y$ 满足与上题完全相同的齐次方程，只是初值由 $CU$ 换成 $C(U-\mathcal E)$。因此
\begin{equation}
  \boxed{I(t)=\frac{\mathcal E-U}{L}\,
  \frac{\e^{s_+t}-\e^{s_-t}}{s_+-s_-}},
\end{equation}
其中 $s_\pm$ 与上题相同。欠阻尼时
\begin{equation}
  \boxed{I(t)=\frac{\mathcal E-U}{L\omega_d}
  \e^{-\alpha t}\sin\omega_dt}.
\end{equation}
若 $U=\mathcal E$，初始电容电压已与电源平衡，回路不产生电流；若 $U>\mathcal E$，上式给出负电流，即电容反向向电源一侧放电。
其中 $U$ 按图中极性定义为有符号初电压。临界阻尼时两个根重合，不能直接代入上面的商式；
利用 $\lim_{s_+\to s_-}(\e^{s_+t}-\e^{s_-t})/(s_+-s_-)=t\e^{s_-t}$，得到
\begin{equation}
 \boxed{I(t)=\frac{\mathcal E-U}{L}t\e^{-Rt/(2L)}},\qquad R=2\sqrt{L/C}.
\end{equation}
过阻尼时用两个不同实根的通式。各分支均满足 $I(0)=0$、$\dot I(0)=(\mathcal E-U)/L$。

图~\ref{ch24:fig:review-rlcdiagram-2}中，图示开关已闭合，C为等长双极板，R、L使用独立元件符号；电流正向与初始电压极性共同决定方程。

\end{xesolution}

\Needspace{6\baselineskip}
\begin{origexamplebox}{例题 9.21【磁悬浮列车】}
把列车等效为一个长为 $a$、宽为 $L$、总电阻为 $R_c$ 且自感可忽略的闭合导体框，质量为 $m$。导轨中存在行波磁场
\begin{equation*}
  B(x,t)=B_0\cos(kx-\omega t),
\end{equation*}
方向垂直导体框平面。

(1) 若某时刻车体左边坐标为 $x$、向右速度为 $v$，求瞬时电磁动力；

(2) 解释这种行波磁场为何能驱动列车。
% Question contract: given frame, speed and field law; no phase-speed result or zero-field construction.
\begin{center}\begin{tikzpicture}[scale=.8,>=Latex]
 \draw[Ink,thick](-4,1.4)--(4,1.4) (-4,-1.4)--(4,-1.4);
 \draw[DeepBlue,very thick](-1.2,-1.15)rectangle(1.2,1.15);
 \node at(0,.25){$m,\ R_c$};
 \draw[vector](1.45,0)--(3,0)node[right]{$v$};
 \draw[<->](-1.2,-1.9)--(1.2,-1.9)node[midway,below]{$a$};
 \draw[<->](-4.4,-1.15)--(-4.4,1.15)node[midway,left]{$L$};
 \node at(0,2.1){$B(x,t)=B_0\cos(kx-\omega t)$，垂直框面};
\end{tikzpicture}\end{center}
\figcap{闭合导体框与绝缘机械支承；框长、框宽及瞬时速度为已知量。}\label{fig:em24-question-21}

\end{origexamplebox}

\begin{xesolution}
\par\noindent\begin{minipage}{\linewidth}
\travelingFieldTrainDiagram
\figcap{行波磁场与绝缘机械支承}\label{ch24:fig:review-travelingfieldtraindiagram-1}
\end{minipage}\par\medskip

令
\begin{equation}
  \xi=k\left(x+\frac a2\right)-\omega t,
  \qquad
  \delta=\frac{ka}{2},
  \qquad
  v_\phi=\frac\omega k.
\end{equation}
穿过导体框的磁通为
\begin{align}
  \Phi(x,t)
  &=L\int_x^{x+a}B_0\cos(kx'-\omega t)\dd x'\\
  &=\frac{2LB_0}{k}\sin\delta\cos\xi.
\end{align}
跟随车体求时间导数，得总感应电动势
\begin{equation}
  \mathcal E=-\frac{\dd\Phi}{\dd t}
  =2LB_0\sin\delta\sin\xi\,(v-v_\phi).
\end{equation}
忽略自感时 $I=\mathcal E/R_c$。对回路做微小平移，磁场力可写为 $F_x=I\,\partial\Phi/\partial x$，而
\begin{equation}
  \frac{\partial\Phi}{\partial x}=-2LB_0\sin\delta\sin\xi.
\end{equation}
因此
\begin{equation}
  \boxed{F_x=\frac{4L^2B_0^2}{R_c}\sin^2\!\delta\,\sin^2\!\xi\,(v_\phi-v)}.
\end{equation}
若速度在一次相位循环中变化很小，可将其近似看成常量。车体看到的相位变化率是
$\dot\xi=kv-\omega$，应在滑差周期 $T_s=2\pi/|\omega-kv|$ 上取平均，而非固定位置的周期 $2\pi/\omega$。
当 $v\ne v_\phi$ 时，在此前提下 $\langle\sin^2\xi\rangle=1/2$，故
\begin{equation}
  \boxed{\langle F_x\rangle=\frac{2L^2B_0^2}{R_c}
  \sin^2\!\frac{ka}{2}\,(v_\phi-v)}.
\end{equation}
当列车慢于行波相速 $v_\phi$ 且耦合非零时，平均力向右；快于行波时平均力反向。
若 $\sin(ka/2)=0$，框长恰跨整数个波长，总磁通恒为零，驱动力对任意速度均为零。
同步时 $v=v_\phi$，磁通也不随车体改变，直接由瞬时式得到零电动势和零力，无须使用发散的滑差周期。

若保留自感 $L_c$，令 $\Omega=\omega-kv$、$K=2LB_0\sin(ka/2)$，在恒速周期稳态下
$L_c\dot I+R_cI=-K\Omega\sin\xi/k$，$F_x=-KI\sin\xi$。
正弦激励的同相响应被因子 $R_c/(R_c^2+\Omega^2L_c^2)$ 缩减，因而
\[
 \langle F_x\rangle=
 \frac{K^2R_c}{2(R_c^2+\Omega^2L_c^2)}(v_\phi-v).
\]
在非零耦合、$R_c>0$ 的稳态平均意义下，方向仍趋向同步，但不能把这个结论套用于任意瞬态。
图中的导轨只作绝缘机械支承，不与闭合导体框构成额外电路；正负磁场区随相位向右移动。
本模型解释电磁推进，不负责解释列车的竖直悬浮力。

图~\ref{ch24:fig:review-travelingfieldtraindiagram-1}中，点、叉区表示某一时刻磁场的相反方向，虚线为零场位置，场图样以相速移动。

\end{xesolution}
